\documentclass[12pt]{article}
\usepackage[utf8]{inputenc}
\usepackage[T1]{fontenc}
\usepackage{hyperref}
\title{Yir-Yoront Language Resources for K-12 Curriculum Development}
\author{Compiled by Claude Sonnet 5.5 (claude-sonnet-5-5) for Andr\'as Kornai}
\date{}
\begin{document}
\maketitle

\section{Introduction}
Yir-Yoront (ISO 639-3: yyr) is a dormant Australian language of the Cape York Peninsula, Queensland. No remaining community members use it as a first language, and it is not known to be taught in schools. The search found no school curriculum materials, primers or community programmes. What exists is scholarly documentation, mostly by Barry Alpher.

\section{Primary Educational Resources}
No primers, textbooks or classroom materials for Yir-Yoront turned up in the search. Any curriculum would have to be built from the scholarly sources below.

\section{Linguistic and Reference Materials}
\begin{itemize}
\item Barry Alpher, \emph{Yir-Yoront Lexicon: Sketch and Dictionary of an Australian Language} (Mouton de Gruyter, 1991). It runs to about 795 pages and covers phonology, orthography, grammar, semantics, totems and kinship. Catalogue record: \url{https://oeaw.ac.at/resources/Record/9783110872651}
\item Another catalogue record for the same book: \url{https://bvbvf1.bib-bvb.de/vufind/Record/DE-604.BV004553057/Details}
\item Alpher 1991, bibliographic entry in Glottolog: \url{https://glottolog.org/resource/reference/id/151393}
\item Alpher 1991, entry in WALS: \url{https://wals.info/refdb/record/Alpher-1991}
\item Barry Alpher, \emph{Son of Ergative: the Yir Yoront Language of Northeast Australia} (Cornell PhD dissertation, 1973), the basis of the 1991 book: \url{https://glottolog.org/resource/reference/id/87539}
\item R. M. W. Dixon's review of the lexicon, \emph{Language} (1995): \url{https://glottolog.org/resource/reference/id/6439}
\item Bruce Rigsby's review of the lexicon, \emph{Canberra Anthropology} (1993): \url{https://glottolog.org/resource/reference/id/6084}
\item A review article on Project MUSE: \url{https://muse.jhu.edu/article/452948/pdf}
\item Overview article: \url{https://en.wikipedia.org/wiki/Yir-Yoront_language}
\end{itemize}

\section{Religious and Community Resources}
The search found no scripture translations or community organisations for Yir-Yoront.

\section{Digital and Audio Resources}
Search results indicate that Paradisec holds Yir Yoront recordings in the Barry Alpher collection (BA1) and the Arthur Capell collection (AC1). I did not see a direct Paradisec URL, so none is given here. The search also indicated that AIATSIS lists the 1973 dissertation as a closed-access manuscript.

\section{Curriculum Development Context}
The language is dormant and no first-language speakers remain, so a school curriculum would be a revitalisation effort. It would depend on Alpher's lexicon and sketch grammar and on archived recordings. Descendant-community involvement and permissions would be needed. Little else was found.
\end{document}